\section*{Глава 1. Анализ предметной области}
\addcontentsline{toc}{section}{Глава 1. Анализ предметной области}
\stepcounter{section}

\subsection{Обзор существующих интернет-магазинов спортивных товаров}

На современном рынке интернет-торговли спортивными товарами можно выделить две группы решений: крупные универсальные маркетплейсы и специализированные розничные магазины. Рассмотрим ключевые особенности трёх известных игроков.

\textbf{Wildberries} – международный маркетплейс, предлагающий товары повседневного спроса, включая спортивную одежду, обувь и инвентарь. Платформа объединяет множество продавцов и обеспечивает централизованную обработку заказов. Личный кабинет пользователя позволяет отслеживать покупки, оформлять возвраты, участвовать в акциях. Для малого бизнеса доступна модель размещения товаров с комиссией, однако полноценная автономная система управления заказами и клиентской базой отсутствует – все процессы жёстко интегрированы в инфраструктуру маркетплейса.

\textbf{Decathlon} – вертикально-интегрированный ритейлер, развивающий собственные спортивные бренды. Интернет-магазин предлагает широкий ассортимент товаров для более чем 70 видов спорта, дополненный программой лояльности и персонализированным личным кабинетом. В нём хранится история заказов, адреса доставки, информация о гарантии. Система управления заказами оптимизирована под розничные продажи одной компании и не тиражируется для сторонних продавцов.

\textbf{SportMaster} – российская розничная сеть, чей онлайн-магазин предоставляет доступ к ассортименту спортивных товаров, бонусную программу и личный кабинет с отслеживанием статуса заказа. Платформа обслуживает исключительно собственные операции сети; инструменты для создания независимого магазина или адаптации под нужды небольших специализированных продавцов отсутствуют.

Сравнительная характеристика рассмотренных решений по ключевым критериям приведена в таблице~\ref{tab:competitors}.

\begin{table}[H]
    \centering
    \caption{Сравнительный анализ интернет-магазинов спортивных товаров}
    \label{tab:competitors}
    \begin{tabular}{|p{2.3cm}|p{2.8cm}|p{3.3cm}|p{2.9cm}|p{3.2cm}|}
        \hline
        \textbf{Название} & \textbf{Ассортимент} & \textbf{Управление заказами} & \textbf{Личный кабинет} & \textbf{Гибкость для малого бизнеса} \\ \hline
        Wildberries & Товары всех категорий, включая спорт & Комиссия для продавцов & История заказов, возвраты, избранное & Зависимость от правил маркетплейса, нет автономии \\ \hline
        Decathlon & Спорт и активный отдых, собственные бренды & Сквозной цикл «заказ-доставка», самовывоз & Персональные рекомендации, гарантия, лояльность & Не предусмотрена (система только для Decathlon) \\ \hline
        SportMaster & Спортивная одежда, обувь, инвентарь & Обработка заказов только для своих продаж & Бонусный счёт, история, статус заказа & Тиражирование для сторонних продавцов отсутствует \\ \hline
    \end{tabular}
\end{table}

Анализ показывает, что все рассмотренные платформы либо являются универсальными 
маркетплейсами, либо представляют собой закрытые системы крупных розничных сетей. 
Ни одно из решений не предоставляет малому специализированному магазину спортивных 
товаров готового инструмента с собственной системой управления заказами, 
полноценным личным кабинетом пользователя и возможностью гибкой настройки под 
конкретный ассортимент. Это подтверждает необходимость разработки отдельного 
веб-приложения, ориентированного на небольшие спортивные магазины и сочетающего 
функции управления заказами с персональным клиентским сервисом.






\subsection{Анализ программных инструментов разработки}

Для реализации проекта рассматриваются современные полнофункциональные (Fullstack)
подходы. Выбор инструментов обусловлен требованиями к производительности и
масштабируемости:




\textbf{Фронтенд-фреймворк и мета-фреймворк.} При выборе инструмента для построения пользовательского интерфейса сравнивались React, Vue.js и Angular (таблица~\ref{tab:frontend_compare}).

\begin{table}[H]
    \centering
    \caption{Сравнение фронтенд-фреймворков}
    \label{tab:frontend_compare}
    \small
    \begin{tabular}{|p{2.5cm}|p{4cm}|p{4cm}|p{4cm}|}
        \hline
        \textbf{Критерий} & \textbf{React (выбран)} & \textbf{Vue.js} & \textbf{Angular} \\ \hline
        Экосистема библиотек для e-commerce & Широкая: React Hook Form, Zustand, React Query, Next.js Commerce & Ограниченная; многие решения не имеют типовой поддержки & Меньше специализированных решений, используется Angular Material с доработками \\ \hline
        Производите- \newline льность и SEO & Next.js с серверным рендерингом (SSR) и статической генерацией (SSG) & Nuxt.js (SSR/SSG), но меньшая популярность для крупных магазинов & Angular Universal для SSR, но настройка сложнее и ресурсоёмкость выше \\ \hline
        Кривая обучения и сообщество & Умеренная; огромное сообщество, множество готовых компонентов & Низкая; простота старта, но меньше специалистов & Высокая; требует знания TypeScript, RxJS, DI \\ \hline
    \end{tabular}
\end{table}

React совместно с Next.js обеспечивает серверный рендеринг страниц каталога товаров и маршрутизацию разделов личного кабинета и системы управления заказами без дополнительных надстроек. Широкая экосистема библиотек (React Hook Form для форм оформления заказа, Zustand для управления состоянием корзины) позволяет ускорить разработку интерфейсов электронной коммерции~\cite{react_docs}.

\textbf{Бэкенд-платформа.} Сравнивались язык Go, Node.js (Express/NestJS) и Python (Django/Flask) по критериям производительности и удобства обработки конкурентных запросов.

\begin{itemize}
    \item \textbf{Go (выбран)}: компилируемый язык со встроенной моделью конкурентности (горутины, каналы). Позволяет эффективно обрабатывать тысячи одновременных запросов при минимальном потреблении памяти, что критично для системы управления заказами и личных кабинетов в периоды пиковых нагрузок~\cite{donovan2015go}. Статическая типизация и быстрая сборка снижают количество ошибок на этапе выполнения.
    \item \textbf{Node.js}: событийно-ориентированная модель с асинхронным I/O, но однопоточный цикл обработки событий требует аккуратного управления длительными операциями; масштабирование на многоядерных системах реализуется через кластеризацию, что усложняет архитектуру.
    \item \textbf{Python}: удобен для прототипирования и работы с данными, однако производительность стандартных веб-фреймворков (Django) уступает Go при высоких нагрузках, а асинхронные возможности (FastAPI) всё ещё требуют тщательной настройки.
\end{itemize}

Таким образом, Go обеспечивает наилучший баланс между производительностью, простотой конкурентного программирования и удобством развёртывания конечного приложения.

\textbf{Формат обмена данными.} Для взаимодействия клиентской и серверной частей выбран формат JSON (JavaScript Object Notation). В сравнении с XML он обладает меньшей избыточностью, легко читается как человеком, так и машиной, а встроенная поддержка в JavaScript и Go минимизирует накладные расходы на сериализацию/десериализацию. Это особенно важно при передаче структурированных данных о товарах, заказах и профилях пользователей через REST API. Альтернативы, такие как Protocol Buffers (protobuf), обеспечивают более компактное представление, но требуют дополнительных этапов кодогенерации и менее удобны на этапе отладки.

Представленный стек (Next.js/React + Go + JSON) формирует целостную архитектуру, в которой каждый компонент выбран по совокупности критериев: производительности, развитости экосистемы, простоте интеграции и поддержки со стороны сообщества разработчиков.
















% \begin{itemize}
% 	\item \textbf{Frontend-фреймворк:} Сравнение React и Vue.js показало преимущество React в 
% 	части экосистемы библиотек для разработки интерфейсов электронной коммерции, управления 
% 	состоянием приложения и работы с формами (напр., \texttt{React Hook Form, Zustand}). 
% 	В качестве основного фреймворка выбран React совместно с Next.js, обеспечивающим 
% 	серверный рендеринг страниц каталога товаров и маршрутизацию между разделами личного 
% 	кабинета и системы управления заказами.
% 	\item \textbf{Backend-платформа:} Использование языка \textit{Go} позволяет
% 	организовать высокопроизводительную асинхронную обработку HTTP-запросов
% 	благодаря встроенной модели горутин, что критично при одновременной обработке
% 	множества заказов и запросов к личным кабинетам пользователей.
% 	\item \textbf{Формат обмена данными:} Использование формата \texttt{JSON} для передачи
% 	данных о товарах, заказах и профилях пользователей между сервером и клиентом
% 	обеспечивает простоту интеграции и читаемость структуры REST API.
% \end{itemize}

% Для реализации клиентской части приложения выбран фреймворк Next.js на основе библиотеки React. 
% В процессе настройки компонентов каталога товаров, страниц личного кабинета и модулей управления 
% заказами использовались официальные руководства по быстрому старту из документации разработчиков 
% React и Next.js.~\cite{react_docs}.

% Для реализации серверной части приложения выбран язык \textit{Go},
% архитектура которого подробно описана в работе~\cite{donovan2015go}.
% Использование встроенной модели конкурентного выполнения на основе горутин
% позволяет эффективно обрабатывать запросы от множества пользователей
% одновременно, что особенно актуально при высокой нагрузке в периоды распродаж.

\subsection{Формулировка цели и задач работы}

На основе проведенного анализа инструментов (см. подраздел 1.2) и выявленных
пробелов в существующих продуктах, была сформулирована цель исследования.

\textbf{Цель работы} --- создание веб-приложения интернет-магазина спортивных
товаров, обеспечивающего управление каталогом продукции, обработку заказов
и предоставление пользователям персональных личных кабинетов.


Для достижения поставленной цели необходимо решить следующие \textbf{задачи}:
\begin{enumerate}
	\item Провести сравнительный анализ существующих интернет-магазинов
          спортивных товаров и обзор современных технологий разработки
          веб-приложений;
    \item Разработать архитектуру клиент-серверного приложения и спроектировать
          структуру базы данных для хранения товаров, заказов и данных
          пользователей;
    \item Реализовать программные модули каталога товаров, корзины и системы
          управления заказами;
    \item Реализовать систему личных кабинетов пользователей с возможностью
          просмотра истории заказов и редактирования профиля;
    \item Провести тестирование разработанного веб-приложения и оценить
          корректность работы основных функциональных модулей.
\end{enumerate}